\chapter{Discussion}
\label{ch:discussion}

Chapter~\ref{ch:evaluation} reports what was measured. This chapter interprets
it: what the six verdicts mean taken together, how far each estimate can be
pushed, and what the study cannot support. Section~\ref{sec:discussion_core}
states the central finding and Section~\ref{sec:discussion_h1h6} resolves the
tension the pre-registration deliberately built between H1 and H6.
Sections~\ref{sec:external_interpretability} to~\ref{sec:ablation_establishes}
bound the three analyses whose readings are most easily over-extended, the
external estimate, the anchor-independent label, and the feature ablation.
Section~\ref{sec:implications} draws out what the findings imply for how
studies of this kind should be reported, and
Section~\ref{sec:limitations} sets out the limitations in full.


\section{The Core Finding}
\label{sec:discussion_core}

The most important finding is definitional rather than numerical. Under a
conjunction-based onset rule the Sepsis-3 label on MIMIC-IV \emph{cannot}
contain an event before clinical action begins (Proposition~\ref{prop:h3}), and
under the standard $\min(t_{\text{susp}}, t_{\text{SOFA}})$ rule applied to the
same stays, whose septic membership is cohort-wide identical, it contains
\pcPreTreatOnsetsVarF{} of them, separable from
their surrounding at-risk hours at \pcPreTreatAurocPrimaryF{} AUROC on the
onset-hour estimand. A published MIMIC-IV sepsis result is
therefore only as informative about physiological deterioration as its onset
rule permits, and a result reported without the timing branch is measuring
discrimination around a clinician-action-defined onset, inside a window that
action opened.
The Utility Score stays close to the no-alert strategy for every model, and
one reason for that is a property of the score itself rather than of the data:
it charges for every false-alert hour while crediting only the first correct
one, so at an hourly event rate below 0.2 per cent the false-alert term
dominates any attainable reward. Whether the temporal split contributes a
second reason, by exposing the models to an admission era in which
treatment-related signals are recorded differently, is a plausible explanation
and not one this study tests: no analysis here decomposes the score into
era-related and metric-related components. AUROC survives the same conditions
because the ordinal ranking of patients is preserved even as the achievable
alert volume becomes unusable, and that contrast is what the metric-choice
finding rests on.

That divergence is what Wang et al.\ \cite{wang2025srpm} documented across
ninety-one models moving from internal to external validation, and what
Section~\ref{sec:auroc_illusion} identified as the field's central evaluation
failure. The result here locates it one step earlier. On MIMIC-IV, under an
honest time-based split, discrimination and clinical value have already parted
company \emph{internally}: there was scarcely any margin left for external
validation to destroy.


\section{H1 and H6 Together}
\label{sec:discussion_h1h6}

H1 (label dominance) and H6 (model-class null) are in productive tension, as
designed, and the resolution is sharper than either hypothesis anticipated,
because both resolve in the same direction, which the pre-registration treated
as unlikely.

H6 is \emph{not} rejected: the two flexible estimators differ by
\pcHSixAbsEst{} AUROC, and the interval's upper limit sits inside the 0.02
bound, so the boosted trees are not better by more than the margin. H1 is also not
confirmed, and by a wide margin in the direction opposite to the one
hypothesised: the model-class spread \pcSpreadModelClass{} exceeds the label
spread \pcSpreadLabel{} by \pcHOneEst, with an interval
[\pcHOneCiLo, \pcHOneCiHi] that excludes zero.

Read together these are not in conflict; they locate the model-class effect
precisely. The large model-class spread is the distance between \emph{families}
of estimator, a six-item bedside score against a fitted hazard model, not
between one fitted model and another. Between the two flexible estimators
evaluated here the difference is small in absolute terms, though the paired
contrast favours the primary model and two-sided equivalence is not established
(H6); across families it is the single largest factor in the decomposition
(H1). The conventional literature, which fixes the label and
compares flexible models to each other, is therefore measuring the one contrast
that has almost no room left in it, while the choice that would move the number (fit anything flexible at all, rather than deploy a score) is not in
dispute and the choice that moves the cohort, the label, is invisible to
the metric being reported.

The H6 result also settles, for this task and this dataset, the narrower
version of the statistics-versus-machine-learning question left open in
Section~\ref{sec:false_dichotomy}. Christodoulou et al.'s
\cite{christodoulou2019cox} equivalence finding was established for binary
outcomes on tabular data, and Fagerstr{\"o}m et al.\ \cite{fagerstrom2019lisep}
showed it breaking on the temporal early-warning axis, where a deep model
detected patients up to twenty hours earlier than a Cox baseline. Here, with
the label, the split, the feature set and the evaluation held fixed, an
interpretable discrete-time hazard model is not outperformed by
gradient-boosted trees, and is ahead of them by \pcHSixAbsEst{} AUROC on an
interval that excludes zero. That is not
a contradiction of Fagerstr{\"o}m et al.: their comparison was against a static
Cox specification, which is exactly the comparator this study retains as the
deliberately weak one, and which loses by a wide margin here too
(\pcAurocCoxB{} against \pcAurocPrimaryB). The result holds between the primary
model and the boosted trees; nothing of the kind holds between either of them
and the field's default survival specification.

This is also the clearest illustration in the thesis of what the inference
layer adds and what it does not. On point estimates alone H1 was already
descriptively refuted, and refuted in the direction opposite to the one the
study set out expecting: the label spread is \pcSpreadLabel{} against a
model-class spread of \pcSpreadModelClass. What the interval establishes is
that the reversal is not an artefact of sampling: the label spread's interval
is [\pcSpreadLabelCiLo, \pcSpreadLabelCiHi], the model-class spread's is
[\pcSpreadModelClassCiLo, \pcSpreadModelClassCiHi], and the two do not
overlap. The two sources are not of comparable size, and the defensible
finding is the less convenient one: the label moves discrimination by roughly
an order of magnitude less than the model class, while moving the cohort, the
event count and the denominator of every clinical quantity computed on it. The
label effect is real and consequential, but it is not an effect on
learnability, which is the only channel through which the field's usual metric
could have detected it.


\section{What the External Estimate Can Support}
\label{sec:external_interpretability}

Discrimination is estimated on \pcExtNEventsA{}, \pcExtNEventsB{} and
\pcExtNEventsC{} external events under Variants A, B and C
(Section~\ref{sec:external_validation}). The interpretability floor adopted
here is \pcExtMinEvents{} events, following the external-validation sample-size
guidance of Riley et al.\ \cite{riley2019sample} and Collins et al.\
\cite{collins2024tripodai}: below it the confidence interval on an external
AUROC is wider than the degradation the study is trying to detect.
\textbf{All three variants fall below it.} The degradations are accordingly
stated as directions, not as magnitudes, and the H5 interval is the formal
expression of that limit.

One external estimate in this study does clear the floor, and it is important
not to mistake it for a Sepsis-3 result. The anchor-independent arm of
Section~\ref{sec:external_altlabel} scores \pcExtAltNEvents{} events, because
its label needs no culture record and therefore reaches the whole cohort. It is
post-hoc, it targets acute organ dysfunction rather than Sepsis-3, and it is
excluded from H5. What its event count buys is precision about the
\emph{direction and rough size} of external degradation, which the three
pre-registered variants cannot supply; it buys nothing about the Sepsis-3
target, for which the floor still binds.

Two further quantities are unavailable at this event count in any case. AUPRC is
dominated by sampling noise at a few dozen positives, and the Utility Score
requires enough septic stays for the timeliness reward to be estimated at all.
Neither is interpreted externally.

\paragraph{The antibiotic-only arm.} Relaxing the culture requirement across
the full cohort raises the event count by two orders of magnitude (\pcExtAbxNEventsA, \pcExtAbxNEventsB{} and \pcExtAbxNEventsC{} events across
\pcExtAbxNRows{} person-hours), and therefore clears the floor with room to
spare. Under that anchor the primary model scores \pcExtAbxAurocPrimaryA,
\pcExtAbxAurocPrimaryB{} and \pcExtAbxAurocPrimaryC.

The comparators in the same arm carry a result the Sepsis-3 arm is too small to
see, and it does not favour the flexible estimator. The boosted model scores
\pcExtAbxAurocGbtA, \pcExtAbxAurocGbtB{} and \pcExtAbxAurocGbtC, close to
chance under the two narrower labels, while NEWS2 scores
\pcExtAbxAurocNewsA, \pcExtAbxAurocNewsB{} and \pcExtAbxAurocNewsC. A bedside
score therefore ranks ahead of the boosted model under every label in this arm.
The internal model-class ordering does not transport. That is worth stating
plainly for two reasons. It qualifies H6 in a way the internal analysis cannot:
the two flexible estimators are close in magnitude inside the development
cohort, with the paired contrast favouring the primary model, but only one of
them stays ahead of a bedside score outside it, and it is the same one. And it cuts against this study's own
comparators as much as for them, since a rule-based score overtaking a fitted
model externally is not the result a modelling thesis would choose. The event
count supports the ordering; it does not license a claim about the size of any
gap, for the same reason the Sepsis-3 arm does not, and the arm still labels a
different target.

This arm labels a broader and less specific target: treated suspicion of
infection, without the organ-dysfunction conjunction's culture limb. An AUROC
computed against it therefore answers a different question and is never
compared like-for-like with the internal Sepsis-3 estimates, nor pooled with the
primary arm. What it does supply is a consistency check that the primary arm
cannot supply for itself. The degradation it shows is of the same sign, the same
label ordering and a somewhat larger magnitude than the Sepsis-3 arm's, on a
sample large enough to resolve it. Two estimates that disagree by two orders of
magnitude in event count and agree on direction are better evidence of a
direction than either alone would be. Because the two arms label different
targets, that agreement is descriptive: it does not establish degradation for
the Sepsis-3 target, and the Sepsis-3 interval covers zero. It is the strongest
statement the external analysis supports, and it is deliberately weaker than a
claim that discrimination degrades or by how much.

Against the published record assembled in Section~\ref{sec:generalisation},
the direction is unsurprising and the magnitude is unresolvable. Moor et al.\
\cite{moor2023review} report an external penalty of 0.085 AUROC across four
international databases; the point estimate here, \pcHFiveEst, is about
three-quarters of that. Where this study cannot follow Moor et al.\ is in
attributing the loss, because their harmonised databases shared labelling
inputs and eICU-CRD does not.


\section{What Variant D Does and Does Not Establish}
\label{sec:variant_d_establishes}

Variant D's high AUROC must be read with the limitation stated in advance in
Section~\ref{sec:alt_labels}. Its label is a threshold function of recorded
physiology, and five of the six SOFA components are themselves model inputs, so
its discrimination is partly self-fulfilling. The \emph{level} of
\pcAltAurocFullPrimaryD{} (\pcAltAurocFullGbtD{} for GBT) is not evidence that
deterioration prediction is a solved problem, and it is not comparable with
Variant B's \pcAurocPrimaryB.

Three things it does establish:

\begin{enumerate}
    \item \textbf{The label sensitivity is not wholly definitional.} At
    $\kappa = \pcKappaVsBD$ and a \pcPctExactMatchD\% onset match, an
    anchor-independent label
    disagrees with Sepsis-3 about both membership and timing. Had D closely
    tracked B, the H1 result would have been an artefact of how one definition
    is read; it is not.

    \item \textbf{The alert burden is label-invariant; the utility ceiling is
    not.} Both alternative labels are markedly more scoreable than Sepsis-3.
    The best achievable Utility Score rises from \pcUtilityMaxGbtB{} under
    Variant B to \pcUtilityMaxGbtE{} under E and \pcUtilityMaxGbtD{} under D,
    with the primary model following the same pattern
    (\pcUtilityMaxPrimaryE{} and \pcUtilityMaxPrimaryD). What does not improve
    to a deployable level is the alarm volume: even D's best configuration
    carries \pcUtilityOptBurdenGbtD{} false alerts per true alert, better than
    Variant B's \pcUtilityOptBurdenGbtB{} but still nine times the 1.4
    benchmark of Moor et al.\ \cite{moor2023review}. Contribution C3 survives
    the change of label in its alarm-volume form, which the pre-registered
    variants alone could not show, and does not survive in the form of a
    utility ceiling near zero. The rise is itself expected from the caveat
    above: a label built from recorded physiology is easier to alert on than
    one built from a clinician's decision.

    \item \textbf{The two metrics order the models differently, and the
    ordering is stable across labels.} Evaluated at its own best threshold, GBT
    leads on the Utility Score under B, D and E alike
    (\pcUtilityMaxGbtB{} against \pcUtilityMaxPrimaryB, \pcUtilityMaxGbtD{}
    against \pcUtilityMaxPrimaryD, \pcUtilityMaxGbtE{} against
    \pcUtilityMaxPrimaryE). On AUROC the interpretable model leads under B and
    E, and GBT leads only under D, whose label is a threshold function of the
    physiology both models are given (Section~\ref{sec:h6}). Neither ordering
    is reversed by changing the label; what changes the ordering is changing
    the metric. A third and different reversal appears if every model is scored
    at the same 0.5 cut-off, which rewards a model for being too poorly
    calibrated to reach it. The concern H5 was written to address, that the
    verdict depends on the metric reported, is what these three orderings
    illustrate; they are not a test of H5's internal-to-external contrast,
    which is reported in Section~\ref{sec:h5}.
\end{enumerate}

A fourth item was added after the first three, under Protocol Amendment~9, and
it is the only one that speaks to the external database as a whole.

\textbf{The external event-rate gap is largely attributable to the anchor
rather than to the population.} The Sepsis-3 arm reaches only \pcExtNMicroStays{} of
\pcExtAltNStays{} eICU-CRD stays and yields an event rate
\pcExtRateRatioB{} times the development cohort's. Two readings of that gap
have very different consequences. Either eICU's patients differ from
MIMIC-IV's to a degree that would undermine any external comparison, or the
Sepsis-3 label is simply constructible less often there because it requires a
documented culture. Nothing in the Sepsis-3 arm can separate them, since both
readings predict exactly the same shortfall.

Variant D separates them, because it holds the physiology and the two databases
fixed and removes only the treatment anchor. Applied to all \pcExtAltNStays{}
stays it gives \pcExtAltRateEicu{} events per 1{,}000 person-hours against
\pcExtAltRateMimic{} internally, a ratio of \pcExtAltRateRatio{} where the
anchored label gave \pcExtRateRatioB. Most of the shortfall disappears when the
label stops requiring a record of what a clinician did. This is the direct
evidence for Limitation~\ref{lim:label_transport}'s claim that the binding
external constraint is label transport rather than cohort comparability, and
the pre-registered variants could not have produced it: all three share the
anchor whose contribution is in question.

The design does not isolate the anchor from everything else, and the residual
is where the remaining explanations live. Removing the anchor closes the gap
from a ratio of \pcExtRateRatioB{} to \pcExtAltRateRatio, which leaves roughly
a quarter of the development cohort's rate unaccounted for. Case mix,
measurement and charting frequency, the coarser five-component SOFA delta eICU
supports, and genuine differences in illness severity all sit inside that
residual and none of them is separated here. The defensible statement is
therefore comparative and not exclusive: the anchor accounts for most of the
gap. The warrant for reading it that way is that the two labels are applied to
the \emph{same} two cohorts, so a population difference large enough to produce
the anchored ratio of \pcExtRateRatioB{} would have had to produce a comparable
shortfall under the anchor-independent label as well, and it does not. That
bounds the anchor's contribution rather than isolating it, and it inherits the
deterioration label's own cross-database difference: the two constructions are
not strictly identical, so the comparison is between two labels that each
travel differently, not between one label and a fixed yardstick.

Its discrimination must be read under the same restriction as the rest of this
section. Variant D's external AUROC of \pcExtAltAurocPrimary{}
(\pcExtAltAurocGbt{} for GBT) is measured on \pcExtAltNEvents{} events rather
than \pcExtNEventsB, so it is the one external estimate here that is not
limited by its event count. That buys precision about a deterioration rule, not
about sepsis. The label remains a threshold on physiology the models are given,
so the \emph{level} is no more interpretable as sepsis performance than it was
internally, and it is not comparable with \pcExtAurocPrimaryB.

What the event count does support is the \emph{drop}, and under Protocol
Amendment 10 it supports it with an interval. The primary model loses
\pcExtAltDropPrimary{} of AUROC on transport, 95\% CI
\pcExtAltDropCiPrimary, and the boosted comparator \pcExtAltDropGbt{}
(\pcExtAltDropCiGbt); both intervals exclude zero and both are narrower than a
hundredth of a point. This is the study's one demonstrated transport loss, and
its scope is narrow in a way worth stating twice: it is demonstrated for a
deterioration rule, it says nothing about Sepsis-3, and it cannot be pooled
with or compared against the anchored arm. Its use is negative. It removes one
reading of the Sepsis-3 arm's inconclusiveness, namely that these two databases
are similar enough for a frozen model to move between them at no cost, and
leaves the reading the event count already implied, that
\pcExtNEventsB{} events cannot resolve a degradation of the size the
literature anticipates. On the target that has the events, the degradation is
there.


\section{What the Ablation Establishes}
\label{sec:ablation_establishes}

The ablation is a small experiment and should not be asked to carry more than
it can. It is exploratory, it is run on one cohort, and with \pcHTwoNTerms{}
covariates the stability counts on either side are small enough that the
comparison is descriptive rather than tested.

What it does establish is a distinction the thesis needs and could not
otherwise make. The argument that the sepsis label is constituted by clinician
action invites an obvious objection: that the same is true of everything in an
ICU dataset, so the observation is unfalsifiable and therefore empty. The
ablation answers that objection empirically. The feature set divides into
predictors that record the patient and predictors that record the clinician,
the division is made in advance rather than fitted, and the two behave
differently under a change of label. That is a testable version of the claim,
and it is tested here rather than asserted.

It could have failed, and the form of the result matters as much as its
direction. Had discrimination fallen substantially, the finding would have been
that the models lean on clinician behaviour: supportive of the same argument,
but a concern about what reported performance is measuring, belonging as much in
the limitations as in the results. Had stability been unchanged, the
concentration of instability in the action-derived features would have been
small-sample noise and the reading would have had to be withdrawn. Neither
happened. Discrimination falls by \pcAblDeltaAurocPrimaryB{} for the primary
model under the primary label, and by less than 0.01 in all six comparisons,
while the unstable count falls from \pcHTwoEst{} to \pcAblHTwoNUnstable. The
asymmetry between those two movements is what separates the two hypotheses:
features carrying the model's predictive weight could not be removed at a cost
this small, and features whose removal takes five of six unstable terms with
them are where the label sensitivity lives.

The reading this licenses is bounded and worth stating exactly. The
discrimination cost is small but it is not zero, and five of the six paired
intervals exclude zero
(Section~\ref{sec:ablation}), so this is a statement about magnitude and not
about absence. It is \emph{not}
that the action-derived features are unimportant: a feature that
moves AUROC by less than 0.01 may still be the reason a particular alert fires.
It is that on this cohort the discrimination such a model reports depends on
them far less than its cross-label stability does. A published effect
estimate for a treatment indicator from a model of this kind is therefore
reporting something about the labelling regime as much as about the patient, and
the cheapest available remedy, dropping the explicit action features, moves
AUROC by \pcAblDeltaAurocPrimaryC{} in the worst of the six paired comparisons
and by less in the other five, while removing most of the cross-label
coefficient instability
(\pcHTwoPctUnstable\% of terms flagged in the full fit against
\pcAblHTwoPctUnstable\% in the ablated one). That is a small cost, not an
absent one: \pcAblNExcludeZero{} of the \pcAblNContrasts{} contrasts exclude
zero, and no margin of practical equivalence was pre-registered against which
to call the loss negligible.

The corresponding limitation is stated in Section~\ref{sec:limitations}: the
retained laboratory values are themselves forward-filled from tests somebody
chose to order, so the boundary drawn here is conservative rather than clean. A
model with no trace of clinician behaviour in it is not constructible from this
data, and the ablation measures the effect of removing the explicit traces, not
all of them.


\section{Implications for How These Studies Are Reported}
\label{sec:implications}

The preceding sections bound what this study establishes. This one states what
follows for the design of the next one, because the findings bear less on which
model to build than on what a sepsis prediction study has to report before its
headline number can be read at all.

\paragraph{A label variant is not a sensitivity analysis.}
Section~\ref{sec:label_object} observed that the literature demonstrates
\emph{that} the label matters without attaching an interval to how much. The
result here sharpens the reason that gap is costly. Label choice moves AUROC by
\pcSpreadLabel{} and cohort membership by far more, the two narrowest
readings agree at $\kappa = \pcKappaVsBA$ at near-identical prevalence, so a
study that reports one definition's AUROC has reported a number that would
barely move under a different definition while every count, rate and
denominator beneath it moved substantially. The consequence is that reporting a
second variant as a robustness check answers the wrong question. What needs
reporting is the agreement between the variants, because that is where the
variation is, and it is invisible in the metric the field publishes.

\paragraph{An onset rule is a claim about what can be predicted, and belongs in
the methods.}
Section~\ref{sec:label_action} recorded Kamran et al.'s finding that models
partly learn to predict clinician behaviour, and this study locates the
mechanism one level down: not in the feature set but in the timestamp. Two
studies can share a cohort, a definition and a model and still differ on
whether pre-onset prediction is measurable, purely through the choice between
$t_{\text{susp}}$ and $\min(t_{\text{susp}}, t_{\text{SOFA}})$. That choice is
routinely left implicit in a sentence about Sepsis-3. It should be stated
explicitly, with its consequence for the evaluation window stated alongside it,
because under a conjunction rule an empty pre-treatment window is not evidence
of anything (Proposition~\ref{prop:h3}).

\paragraph{Discrimination and clinical value must both be reported, and neither
substitutes.}
Section~\ref{sec:auroc_illusion} set out why AUROC conceals the failure Wang et
al.\ documented: it is monotone-invariant and dominated by true negatives.
This study adds the observation that a single threshold conceals it a second
way. At the event rate here the conventional 0.5 cut-off separates models by
calibration rather than by clinical value, and the primary model's zero at that
cut-off is an artefact of never alerting there. A threshold-dependent metric is
interpretable only as a swept ceiling with the attaining operating point
reported beside it, and the alert burden at that point (\pcUtilityOptBurdenGbtB{} false alerts per true one here, against the 1.4
benchmark of Moor et al.\ \cite{moor2023review}) is the number that decides
deployability. It appears in almost none of the ninety-one models Wang et al.\
\cite{wang2025srpm} reviewed.

\paragraph{An external cohort should be characterised before it is scored.}
Section~\ref{sec:generalisation} assembled the published external penalties
into a single frame and treated them as comparable. The experience here is that
they are not automatically so. The binding constraint on the external estimate
in this study is not sample size but whether the label's inputs exist in the
destination database at all, and that constraint is invisible in a reported
AUROC and in a reported cohort size alike. A study transporting a
treatment-anchored label should report the coverage of each labelling input in
the destination cohort before it reports any performance number, and should
state the event count the estimate rests on rather than the stay count.

\paragraph{Model-class comparisons are the least informative axis available.}
Section~\ref{sec:false_dichotomy} argued that the statistics-versus-machine-learning
partition does not carve the field. The decomposition supports a stronger
version, bounded by what was actually fitted. Between the two flexible
estimators evaluated here the choice moves AUROC by \pcHSixAbsEst{}, against
\pcSpreadModelClass{} across the six specifications, so on this comparison the
model-class axis has much less room left in it than the range suggests, while
the label choice that sets the cohort goes unreported. Two estimators are a
narrow base for a general claim: no deep temporal model was fitted, and the
boosted comparator was not tuned. The reporting recommendation is
correspondingly bounded: report which flexible estimator was fitted rather than
defending the choice at length, and spend the space on the label and the onset
rule instead.

\paragraph{A subgroup analysis needs an interpretability floor declared in
advance.}
Section~\ref{sec:equity_lit} noted that the nearest prior work measures
subgroup disparity for mortality among already-septic patients rather than for
onset. Extending it to onset prediction exposed a problem that is structural
rather than particular to this cohort: at a low event rate most demographic
strata carry too few events for a subgroup AUROC to mean anything, and an
interval reaching the chance line is easily read as a null result rather than
as an absent measurement. Declaring the floor before estimation, tabulating the
below-floor strata rather than suppressing them, and excluding them from the
corrected family is what keeps the distinction visible. Without it, a study
underpowered for a fairness question can report one as answered.


\section{Limitations}
\label{sec:limitations}

\begin{enumerate}
    \item \textbf{Single-centre development.} All models were developed on
    MIMIC-IV, which originates from a single academic medical centre (Beth
    Israel Deaconess Medical Center, Boston). Clinical practices, patient
    demographics, and EHR documentation patterns at BIDMC may not represent
    those at other institutions. The equity and label-sensitivity findings in
    particular should not be assumed to transfer to hospitals with different
    patient populations or antibiotic prescribing cultures. The eICU-CRD
    analysis is the only external evidence available here, and the limitation
    below bounds how much it can carry.

    \item \textbf{The base cohort is selected, and in two ways that are not
    neutral.}\label{lim:cohort_selection} Only the first ICU stay of each
    patient is retained, and only stays of at least four hours
    (Section~\ref{sec:cohort}). The first restriction removes repeat
    admissions, which are neither a random subset of ICU care nor independent
    of it: a patient returning to intensive care differs systematically from
    one admitted for the first time, and the models are never asked about them.
    The second removes stays too short to build an hourly panel on, and those
    are not a homogeneous group either. Some are rapid deaths, some are rapid
    transfers, and some are low-acuity observation admissions that were never
    going to deteriorate. The two kinds pull the cohort's acuity in opposite
    directions, so the net effect on the reported estimates is not signed here
    and is not recoverable from the analyses reported. What can be said is that
    the population is not all adult ICU admissions but the subset for which an
    hourly real-time task is well posed, and transportability claims should be
    read against that population rather than against ICU care in general.

    \item \textbf{External validation is bounded by label transport, not by
    cohort size.}\label{lim:label_transport} eICU-CRD contains \pcExtNCohortStays{} unit stays, but the
    Sepsis-3 suspicion anchor requires a microbiology culture and
    \texttt{microLab} documents one for only \pcExtMicroCoverage\% of the
    analysis cohort, and both limbs of the anchor together for
    \pcExtBothLimbCoverage\% (\pcExtNBothLimbStays{} stays). The primary
    external analysis is confined to the culture-covered
    sub-cohort, because in the remaining stays no onset is observable under any
    variant; the resulting \pcExtNEventsB{} events under the primary label sit
    below the \pcExtMinEvents{}-event interpretability floor
    (Section~\ref{sec:external_interpretability}). Adding person-hours from the
    uncovered stays would enlarge the denominator without adding a single
    observable onset, so this is not a limitation additional data of the same
    kind would remove, it would require a destination database whose
    microbiology recording matches MIMIC-IV's. That the constraint is the
    anchor rather than the population is shown directly in
    Section~\ref{sec:variant_d_establishes}: a label carrying no treatment
    timestamp recovers an external event rate of \pcExtAltRateRatio{} times the
    development cohort's, against \pcExtRateRatioB{} under the anchored label,
    on the same two databases.

    A second consequence follows, and it concerns composition rather than
    precision. The sub-cohort is defined by a documentation property, and a
    culture is ordered when infection is already suspected, so the
    \pcExtNMicroStays{} covered stays are not a random sample of the
    \pcExtNCohortStays: they are enriched for the patients a clinician had
    begun to investigate. The external estimate therefore describes transport
    into a documentation-selected sub-population rather than into eICU-CRD as a
    whole. The direction of that selection cannot be signed from these data,
    because the uncovered stays carry no observable onset to compare against,
    so it stands as an unquantified limit on generalisability rather than a
    bias with a known sign. The consequence is specific and
    should not be over-read in either direction: the study establishes that
    every external point estimate is lower than its internal counterpart, and
    the Sepsis-3 arm establishes neither that discrimination degrades nor by
    how much.
    External Utility Score and AUPRC are not interpreted at all. The
    antibiotic-only sensitivity arm is well powered and agrees on direction,
    but it labels a different target and cannot stand in for the Sepsis-3
    comparison.

    \item \textbf{The external degradation is not decomposed into its causes.}
    An AUROC that falls on transport may reflect case-mix differences,
    documentation and measurement-frequency differences, genuine differences in
    how sepsis presents across two hundred hospitals, or residual imperfection
    in the transported label. This study measures the total and attributes none
    of it. The antibiotic-only arm bounds the label's contribution loosely (degradation persists when the culture limb is removed, so it is not solely
    an artefact of culture coverage), but a proper attribution would require
    a destination database with equivalent labelling inputs, which is precisely
    what was unavailable.

    \item \textbf{The NEWS2 comparator is not scored identically in the two
    cohorts.} NEWS2 charges three points for a patient who is not alert and
    none for a patient whose consciousness was not assessed. In MIMIC-IV the
    flag is derived from the recorded Glasgow Coma Scale with an unmeasured
    value treated as ``not alert'', so the small minority of person-hours
    carrying no GCS after forward-fill are charged those three points; in
    eICU-CRD, where the variable is absent from the periodic vitals table
    altogether, an unmeasured value is propagated as unmeasured and charged
    nothing. AUROC is a ranking statistic and the internal effect is confined
    to that minority of rows, but the two cohorts' NEWS2 \emph{thresholds} are
    not on a common scale, so the internal and external specificity and
    alert-burden figures for that comparator should be read within a cohort and
    not across the pair. Nothing in the decomposition, the hypothesis verdicts
    or the primary-model results depends on this: NEWS2 enters the study as a
    deliberately weak reference point, and the internal-versus-external
    contrast that H5 tests is computed on the primary model.

    \item \textbf{The feature set cannot be fully purged of clinician action.}
    The ablation in Section~\ref{sec:ablation} removes the
    \pcAblNDropped{} predictors that record clinician behaviour explicitly, vasopressor exposure and the per-test order indicators, but the
    laboratory values it retains are forward-filled from tests somebody chose
    to order, and even a vital sign is recorded on a schedule that reflects how
    closely a patient is being watched. The boundary drawn is therefore
    conservative rather than clean: what is measured is the effect of removing
    the explicit traces of clinical attention, not of removing all of them. A
    predictor set with no such trace is not constructible from routinely
    collected ICU data, which is itself part of the thesis's argument rather
    than only a limitation of it.

    \item \textbf{Deep survival model not implemented.} A deep survival
    comparator was planned but not built. Fagerstr{\"o}m et al.\
    \cite{fagerstrom2019lisep} report a recurrent architecture detecting
    septic-shock onset up to twenty hours earlier than a Cox baseline on the
    same features and target, so the absence of a deep temporal comparator is
    one of two places where H6's result is incomplete. It establishes
    non-inferiority against gradient-boosted trees, not against every flexible
    estimator. The second gap is inferential rather than architectural: the
    contrast is tested one-sidedly, so what is excluded is the boosted trees
    outperforming the hazard model, and two-sided equivalence between the two is
    not established: the interval's lower limit, \pcHSixCiLo, falls outside
    the 0.02 margin (Section~\ref{sec:h6}). A third qualification is the
    comparator's fitting procedure. Its boosting-round count is now chosen on a
    holdout carved out of the training stays (Protocol Amendment 8,
    Section~\ref{sec:gbt_spec}); under the earlier procedure it was chosen on
    the test set, which flattered the comparator. Correcting it moved this
    contrast in the direction that favours the interpretable model, so the
    result should be read as one the uncorrected fits already supported and the
    correction sharpened, rather than as one the correction produced. Neither
    gap nor the correction affects the structural findings about the label.

    \item \textbf{Limited feature set.} Clinical notes, procedure codes, and
    ventilator settings were not included. Richer features might improve
    discrimination but would not change the treatment-constitutionality finding.

    \item \textbf{Part of the label-variant spread is definitional, not
    empirical.} Variants A--C are three readings of one definition and all
    three fire on the same culture-plus-antibiotic anchor. Sepsis-3 makes onset
    a function of treatment timing, so moving the antibiotic--culture windows
    \emph{must} move the onsets; some of the H1 label spread is guaranteed by
    construction and would appear even if the physiology were identical across
    patients. Section~\ref{sec:label_anchor_results} separates the two by
    labelling each limb of the Sepsis-3 conjunction on its own, but the
    separation is post-hoc (Protocol Amendment 2) and rests on two variants,
    not on a systematic survey of alternative definitions.

    \item \label{lim:current_hour}\textbf{Discrimination is measured on the
    current-hour hazard, not on a six-hour horizon.} The models estimate the cause-specific hazard of onset
    in the current hour, and every AUROC, AUPRC and calibration figure in this
    thesis scores that hazard against an indicator marking the onset hour among
    the at-risk hours (Section~\ref{sec:framing}). Three consequences follow.
    The reported discrimination is not comparable, quantity for quantity, with
    published AUROCs from studies that classify a binary ``onset within the next
    $k$ hours'' label, whose positives are $k$ times more numerous and whose task
    is correspondingly easier to score well on; comparisons with such figures in
    Chapter~\ref{ch:review} are therefore indicative rather than exact. The
    covariates entering the prediction for the onset hour are contemporaneous
    with it, so the discrimination figures measure detection at the recorded
    onset rather than anticipation of it, and the quantity that speaks to
    anticipation is the lead-time distribution, which
    Section~\ref{sec:clinical_utility} shows to be uninformative at these alert
    rates. Contemporaneity is exact only to the hour: measurements are binned
    by the elapsed hour in which they were charted, and the onset hour is binned
    the same way, so within that single hour a covariate charted after the onset
    timestamp still enters the prediction assigned to it. Hourly binning is the
    convention in this literature and the effect is confined to the one scored
    hour per septic stay, but it means the onset-hour figures should be read as
    detection at the recorded onset hour and not as evidence of any lead. A
    design that lagged every predictor by one hour, or applied a sub-hour
    timestamp cut-off, would separate the two and is not implemented here. And a cumulative six-hour risk, which is what a deployed instrument
    would display, is a functional of this hazard together with the
    competing-cause hazards and is not reported here. A design that estimated
    and scored both would be stronger; the horizon-free hazard was chosen
    because the competing-risks formulation makes it the more primitive
    quantity, not because the six-hour risk is uninteresting.

    \item \textbf{The primary model's coefficient standard errors are those
    of a penalised estimator, and H2 still carries no uncertainty.} A
    cluster-robust sandwich covariance is computed for the fitted multinomial
    (Section~\ref{sec:primary_model_spec}), but two limits apply. Because the
    specification requires a ridge penalty to be identified at all, it is the
    covariance of the penalised estimator rather than of a maximum-likelihood
    one, so the usual asymptotic interpretation of a Wald interval is
    approximate here. And it bears on coefficients only: H2's estimand is a
    count of covariates crossing a stability threshold across three separate
    fits, which is not a smooth function of one fit's parameters and has no
    analytic variance, so H2 remains a descriptive count with no interval and no
    test. The Amendment 7 diagnostic narrows what can be said about that count
    without removing this limit: it establishes that
    \pcHTwoNNoise{} of the \pcHTwoNFlagged{} flagged covariates move by less
    than their own estimation noise, so the pre-registered count is a loose
    upper bound on the instability present, but the count itself still has no
    sampling distribution and the diagnostic's own ratio is not a test statistic
    (Section~\ref{sec:h2_uncertainty}). Every performance interval in this
    thesis comes from the stay-level cluster bootstrap and is unaffected by
    either limit.

    \item \textbf{The specification requires a penalty to be identified at all.}
    The unpenalised likelihood is completely separated on this design under
    every label variant (Section~\ref{sec:primary_model_spec}), so the reported
    fits are ridge-penalised. The penalty is small and no result depends on its
    value, test AUROC moves by less than 0.02 across a 20-point path, but
    the coefficients are shrunk estimates, not maximum-likelihood ones, and the
    H2 stability counts are properties of the penalised solution. Reporting an
    unpenalised fit here would have been worse, not better: its coefficients
    have no unique maximiser, so cross-label comparison of them is undefined.

    \item \textbf{Feature construction and data cleaning are single-analyst
    work.} The covariate panel is assembled from raw chart and lab events by
    one pipeline, with no independent re-implementation to check it against.
    Two classes of defect are known to be possible in that setting and are
    guarded rather than eliminated: a declared predictor silently failing to
    reach a model because of a name mismatch, and implausible recorded values
    entering the fit. The pipeline now errors on any declared-but-absent
    feature rather than dropping it, and applies an explicit plausibility range
    to all nineteen physiological variables before forward-fill: MIMIC-IV
    chartevents contains values several orders of magnitude outside the
    physiological range, and such rows carry unbounded leverage. Neither guard
    is a substitute for independent replication of the feature pipeline.

    \item \textbf{The pre-registered onset rule is not the standard one.}
    Variants A--C place the onset at $t_{\text{susp}}$ and use the SOFA
    criterion only to gate membership, whereas Seymour et al.\
    \cite{seymour2016assessment} and the PhysioNet/CinC 2019 Challenge
    \cite{reyna2020utility} place it at
    $\min(t_{\text{susp}}, t_{\text{SOFA}})$. The pre-registered family
    therefore cannot address any question about onset \emph{timing} on its own,
    H3 among them: under its rule an empty pre-treatment window is a theorem
    rather than a result (Proposition~\ref{prop:h3}). Variant F supplies the
    measurement the pre-registered design cannot, but it is post-hoc (Protocol
    Amendment 4), it uses the full SOFA score including the vasopressor terms, so its onsets are free of the antibiotic-and-culture anchor but not of
    treatment in general, and its $t_{\text{SOFA}}$ is searched only inside
    the 72-hour modelled panel, so a rise occurring earlier in the admission
    cannot be detected. A design that pre-specified both timing rules across all
    three variants would be stronger than one that adds the second afterwards.

    \item \textbf{The anchor-independent label is not a validated construct.}
    Variant D is an operational deterioration label built for this comparison,
    not a consensus definition with an evidence base behind it, and its onsets
    have no external adjudication. Its label is also a threshold function of
    recorded physiology of which five of six components are model inputs, so
    its discrimination \emph{level} is not comparable with Variant B's, only
    its disagreement with Variant B about which patients deteriorate, and when,
    carries weight.

    \item \textbf{Specific to Sepsis-3 and to MIMIC-IV.} These findings apply
    to the Sepsis-3 definition on MIMIC-IV. SEP-1 and the CDC Adult Sepsis Event
    were not implemented; both are defined on antibiotic days, so neither would
    have tested treatment-independence, but each would give a different cohort
    and might behave differently in other respects. Other databases may behave
    differently again.

    \item \textbf{Operating points are not externally recalibrated.} The
    thresholds in Section~\ref{sec:clinical_utility} are derived on the test
    set at which they are evaluated, so the reported PPV and false-alarm rates
    are optimistic: a prospectively fixed threshold would perform no better.
    The same applies to the maximum Utility Scores in
    Table~\ref{tab:utility_optimum}, which take the best of sixty thresholds on
    the test set and are therefore upper bounds rather than estimates of what a
    deployed system would achieve; they are labelled exploratory for that
    reason. Selecting the threshold on the training years and transferring it
    would be the honest deployment estimate, and would be lower.
    The external operating points in Table~\ref{tab:external_utility} are
    likewise re-derived on eICU rather than transferred from MIMIC-IV, which
    isolates the discrimination loss but hides the additional loss a
    transferred threshold would incur.

    \item \textbf{Alarm episodes are defined by contiguity, not by
    acknowledgement.} False-alarm episodes are counted as contiguous runs of
    alert hours. A real system would apply suppression, snoozing and
    escalation logic, all of which would reduce the episode counts in
    Table~\ref{tab:patient_level}. The alert-hour rates are the conservative
    bound and the episode rates the optimistic one; the true burden lies
    between them.

    \item \textbf{The subgroup discrimination analysis is underpowered, and its
    strata are administrative.} Two distinct limits apply and neither
    substitutes for the other. On precision: only \pcNSubgroupsAboveFloor{} of
    the racial strata carry at least \pcMinSubgroupEvents{} sepsis onsets, the
    floor below which an AUROC is reported but not interpreted here
    (Section~\ref{sec:equity}), and no contrast in family E1 survives
    false-discovery correction. That is a statement about precision, not about
    fairness: the data are consistent both with no disparity and with
    disparities larger than any effect this thesis reports elsewhere. On
    construct validity: MIMIC-IV's \texttt{race} field is a single
    administrative code assigned at registration, so the strata compared are
    documentation categories rather than populations. Broad and granular codes
    drawn from the same population appear as separate levels (``White''
    excludes stays coded ``White -- Other European'', and ``Asian'' excludes
    those coded ``Asian -- Chinese''), so the reference level is a residual
    category and the contrasts are not between disjoint groups. ``Unknown'' and
    ``Unable to obtain'' record the absence of an answer. The
    label-sensitivity result (Table~\ref{tab:equity_label}) is the better
    powered of the two equity analyses because it is a stay-level proportion
    rather than a discrimination metric, but it inherits the same construct
    limitation, and its two significant strata are precisely the two
    missingness categories, which is why it is read in this thesis as a
    finding about documentation completeness rather than about race.

    \item \textbf{The novelty claims rest on a narrative, not systematic,
    search.} The search protocol in Section~\ref{sec:search_strategy} is
    documented and every included source was verified against a primary record,
    but no title-and-abstract screening pass over a full de-duplicated yield was
    performed and no second reviewer screened independently. The ``to the best
    of our knowledge'' claims for contributions C1 and C5 are therefore bounded
    by that coverage: a systematic review could surface a prior study combining
    multi-variant label sensitivity with subgroup equity analysis, and this
    search would not necessarily have found it. Both claims are claims about a
    \emph{combination} of analyses; each component separately has clear
    precedent, which Section~\ref{sec:search_strategy} identifies.

    \item \textbf{Ten protocol amendments are post-hoc.} The inference layer
    (intervals and multiplicity control, Section~\ref{sec:inference}), the
    alternative label limbs D and E (Section~\ref{sec:alt_labels}), the model
    identification and data-plausibility changes
    (Section~\ref{sec:primary_model_spec}), the standard onset timing of
    Variant F (Section~\ref{sec:variant_f}), the action-derived feature
    ablation (Section~\ref{sec:ablation}), the subgroup interpretability
    floor (Section~\ref{sec:equity_strata}), the H2 uncertainty diagnostic
    (Section~\ref{sec:h2_uncertainty}), the comparator fit isolation
    (Section~\ref{sec:gbt_spec}), the full-cohort external arm under Variant D
    (Section~\ref{sec:external_altlabel}) and the intervals attached to three
    quantities previously reported as point estimates
    (Sections~\ref{sec:calibration_results}, \ref{sec:ablation} and
    \ref{sec:external_altlabel}) were all specified after the
    pre-registered analyses had been run and inspected. None changes an
    estimand, threshold, direction or hypothesis, and H1--H6 and families E1/E2
    are computed from unchanged inputs; but a reader should treat the verdicts
    as a pre-specified hypothesis set analysed under a retrospectively specified
    error-control plan, and the amendment results as exploratory. Six of the
    ten warrant a specific note, and they do not all point the same way.
    Amendment 4 enlarges family E3, which tightens rather than relaxes the
    correction applied within it.
    Amendment 6 \emph{reduces} family E1 from thirty contrasts to six, which is
    on its face the manoeuvre the inference plan's own rule against shrinking a
    family prohibits; it is admitted because the criterion is the stratum event
    count alone and is blind to outcome, and because the reduction removed the
    family's only surviving discovery rather than creating one. That is the
    single place where the amendment changes a verdict, and it changes it
    against this study: the discarded contrast rested on twelve events. Nothing
    the reduced family retains survives correction under either denominator, so
    a reader who rejects the reasoning may read the thirty-contrast correction
    instead and will reach the same conclusion about every contrast this
    chapter reports.
    Amendment 7 is the one whose result is \emph{favourable} to the thesis's own
    argument, which is the case for treating it most sceptically. It reports
    that the covariates whose cross-label movement survives comparison with
    their standard errors are the action-derived ones, which is what this thesis
    would like to be true. Three things bound it. It changes no estimand: H2's
    count is reported first and unmodified, and the diagnostic never replaces
    it. It is not a test, for the reasons given in
    Section~\ref{sec:h2_uncertainty}. And it is not independent evidence: it is
    computed from the same three fits as the ablation and shares the same
    flagged set, so it is a second reading of one body of evidence rather than a
    second experiment agreeing with the first.
    Amendment 8 is the only one of the ten that changes a reported estimate
    rather than a family, a diagnostic or a label. It corrects an
    implementation defect: the gradient-boosted comparator's boosting-round
    count was selected by early stopping against the test set, so a
    model-selection choice was being made on the rows the model was then scored
    on. The correction moves the early-stopping set into the training stays,
    split on the stay rather than the person-hour. Its direction was recorded
    before the re-run and is not uniformly convenient. Removing an optimistic
    bias from the comparator widens H6's gap in the interpretable model's
    favour, which is the direction this thesis argues for, so H6 is to be read
    as a finding the uncorrected fits already supported and the correction
    strengthened, not as one the correction produced. It leaves H1 essentially
    untouched, because the model-class range is bounded by the primary model
    above and qSOFA below and the boosted comparator sits at neither extreme.
    Amendment 9 adds an external arm rather than altering one. It applies the
    frozen Variant D models to the whole of eICU-CRD, which is the only way to
    ask whether the external event-rate shortfall follows the anchor or the
    population, and it is registered as descriptive: it joins no corrected
    family, is excluded from H5 on the same terms that exclude Variant D from
    H1, and its AUROC level is not interpretable as sepsis performance for the
    reason given in Section~\ref{sec:variant_d_establishes}.
    Amendment 10 is the one whose direction was fixed before it ran and could
    only run against this thesis. It attaches intervals to three quantities
    previously reported as point estimates and completes a test that had been
    implemented against one boundary of a two-boundary null. An interval placed
    on a point estimate can widen a conclusion and cannot narrow one, and in
    two of the three cases it did: the ablation contrast and one cross-label
    calibration contrast turn out to exclude zero, so the wording that called
    them unaffected has been withdrawn in favour of the magnitudes themselves.
    The amendment log is worth reading with that asymmetry in view, because
    Amendment 7's result flattered this thesis's argument and this one could
    not.

    \item \textbf{Two hypotheses admit no test statistic, and a third is tested
    on a substituted estimand.} H2 is a stability criterion rather than a
    contrast, and the fitted multinomial carries no variance estimate; H3 is
    degenerate because its outcome is absent from the evaluation window
    \emph{by construction} rather than by chance, so no test of it could ever
    have been informative. Both are reported descriptively while still
    consuming a slot in the family size, which makes the correction on the
    remaining four hypotheses stricter, not weaker. H5 is a third and different
    case: its pre-registered form compares two proportional declines, which
    cannot be formed against an internal Utility Score of
    \pcUtilityMaxPrimaryB, so the AUROC limb alone occupies its slot and the
    adjusted $p$-value belongs to that substituted contrast
    (Section~\ref{sec:h5}). The substitution is disclosed in the pre-registration
    rather than reconstructed from the results, but it is a departure from the
    protocol and not merely a degeneracy in it, and a reader is entitled to
    treat H5's verdict as weaker evidence than H1's or H4's for that reason.
    For H1, H4 and H6 the multiplicity correction does not change the
    substantive interpretation: their adjusted $p$-values are \pcHOnePHolm,
    \pcHFourPHolm{} and \pcHSixPHolm, and H1 and H4 are not confirmed from
    their point estimates and intervals, while H6 supports one-sided
    non-inferiority but not two-sided equivalence. The contrast substituted for H5 is the
    exception: its adjusted value is \pcHFivePHolm{} and it is inconclusive
    because its interval is wide and covers zero, though its point estimate is
    in the predicted direction. That is precisely why the substitution needed
    disclosing rather than absorbing.
\end{enumerate}
